% !PRINTED 401
\weberpage{401}
%<<CONT
of the continued fraction development consists of an odd number
of terms. Now a reduced number $\omega = \{a, b, c\}$ is the first
root of
\[
(a, \tfrac{1}{2}b, -c) \quad \text{or} \quad (2a, b, -2c)
\]
(according as $D$ is even or odd) and the second root of
\[
(-a, \tfrac{1}{2}b, c) \quad \text{or} \quad (-2a, b, 2c).
\]
If these two forms are equivalent, that is, if the periods
of $\omega$ consist of an odd number of terms, then to each
two-sided class of forms there corresponds one, and to each pair of
opposite classes two systems of equivalent numbers $\omega$; hence
the class number in the \gesperrt{Gauss}ian sense is just as great as the
number of periods of the reduced numbers $\omega$.

But when the number of terms of the period is even, the
forms $(A, B, C), (-A, B, -C)$ are not equivalent, and the
\gesperrt{Gauss}ian class number is twice as great as the number of
periods of the $\omega$.

\begin{center}
\S.\,129.

\gesperrt{Approximate calculation of the real roots of a
numerical equation by continued fractions.}
\end{center}

On the theory of continued fractions is founded a procedure
due to \gesperrt{Lagrange} for calculating the real roots
of a numerical equation with any desired degree of approximation,
a procedure which can now be set forth in a few words.

Let $f(x) = 0$ be a real equation which has at least one
real root. By one of the various methods of Sections
VIII and IX we can then always find an integer
$a_0$ such that one or more of the roots of $f(x)$ lie between
$a_0$ and $a_0 + 1$. If, then, we transform $f(x)$ by
the substitution
\[
x = a_0 + \frac{1}{x_1}
\]
into $f_1(x_1)$, then $f_1(x_1)$ will have one or more roots which are
greater than 1. We therefore determine a \gesperrt{positive integer}
$a_1$ such that one or more of the roots
of $f_1(x_1)$ lie between $a_1$ and $a_1 + 1$, and transform $f_1(x_1)$ by the
substitution
%CONT>>
